\begin{lemma}
\label{lemma-flat-local-given-residue-field}
Let $(R, \mathfrak m, k)$ be a local ring. Let $K/k$ be a field
extension. There exists a local ring $(R', \mathfrak m', k')$, a flat local
ring map $R \to R'$ such that $\mathfrak m' = \mathfrak mR'$ and such that
$k'$ is isomorphic to $K$ as an extension of $k$.
\end{lemma}

\begin{proof}
Suppose that $k' = k(\alpha)$ is a monogenic extension of $k$.
Then $k'$ is the residue field of a flat local extension $R \subset R'$
as in the lemma. Namely, if $\alpha$ is transcendental over $k$, then we let
$R'$ be the localization of $R[x]$ at the prime $\mathfrak mR[x]$.
If $\alpha$ is algebraic with minimal polynomial
$T^d + \sum \overline{\lambda}_iT^{d - i}$, then we let
$R' = R[T]/(T^d + \sum \lambda_i T^{d - i})$.

\medskip\noindent
Consider the collection of triples $(k', R \to R', \phi)$, where
$k \subset k' \subset K$ is a subfield,
$R \to R'$ is a local ring map as in the lemma, and
$\phi : R' \to k'$ induces an isomorphism $R'/\mathfrak mR' \cong k'$
of $k$-extensions. These form a ``big'' category $\mathcal{C}$ with morphisms
$(k_1, R_1, \phi_1) \to (k_2, R_2, \phi_2)$
given by ring maps $\psi : R_1 \to R_2$ such that
$$
\xymatrix{
R_1 \ar[d]_\psi \ar[r]_{\phi_1} & k_1 \ar[r] & K \ar@{=}[d] \\
R_2 \ar[r]^{\phi_2} & k_2 \ar[r] & K
}
$$
commutes. This implies that $k_1 \subset k_2$.

\medskip\noindent
Suppose that $I$ is a directed set, and
$((R_i, k_i, \phi_i), \psi_{ii'})$ is a system over $I$, see
Categories, Section \ref{categories-section-posets-limits}.
In this case we can consider
$$
R' = \colim_{i \in I} R_i
$$
This is a local ring with maximal ideal $\mathfrak mR'$, and
residue field $k' = \bigcup_{i \in I} k_i$. Moreover, the ring
map $R \to R'$ is flat as it is a colimit of flat maps (and tensor
products commute with directed colimits).
Hence we see that $(R', k', \phi')$ is an ``upper bound'' for the system.

\medskip\noindent
An almost trivial application of Zorn's Lemma would finish the proof
if $\mathcal{C}$ was a set, but it isn't.
(Actually, you can make this work by finding a reasonable bound on the
cardinals of the local rings occurring.)
To get around this problem we choose a well ordering on $K$.
For $x \in K$ we let $K(x)$ be the subfield of $K$ generated
by all elements of $K$ which are $\leq x$.
By transfinite recursion on $x \in K$ we will produce ring maps
$R \subset R(x)$ as in the lemma with residue field extension
$K(x)/k$. Moreover, by construction we will have that
$R(x)$ will contain $R(y)$ for all $y \leq x$.
Namely, if $x$ has a predecessor $x'$, then $K(x) = K(x')[x]$
and hence we can let $R(x') \subset R(x)$ be the local ring extension
constructed in the first paragraph of the proof. If $x$ does not
have a predecessor, then we first set
$R'(x) = \colim_{x' < x} R(x')$ as in the third paragraph
of the proof. The residue field of $R'(x)$ is $K'(x) = \bigcup_{x' < x} K(x')$.
Since $K(x) = K'(x)[x]$ we see that we can use the construction of the
first paragraph of the proof to produce $R'(x) \subset R(x)$.
This finishes the proof of the lemma.
\end{proof}